% GENERATED FILE. Edit canonical lessons, then run npm run generate:books.
% canonical-source-hash: fnv1a64:1e2f32ddf6f4fa3a
% canonical-lessons: FR-C279-rediger, FR-C279-redaction, FR-C279-imprimer, FR-C279-paraitre, FR-C279-parution

\chapter{Écrire pour le journal}
\label{ch:fr-ecrire-pour-le-journal}
\hlchaptermodality{279}

\begin{chapteropening}
\textbf{By the end of this chapter:} \emph{I can say that someone wrote up, printed and brought out a text, and talk about the newsroom and the date it came out.}

Complete the last lesson of chapter 279: I can say that someone wrote up, printed and brought out a text, and talk about the newsroom and the date it came out.
\end{chapteropening}

\section[rédiger]{rédiger — to write, to draft (a text)}

\label{lesson:FR-C279-rediger}



\begin{quote}
Before the new one: say the French for the weekly, then the French for the subscriber.
\end{quote}

\subsection*{You'll want to know: rédiger}
\textbf{rédiger} — \textquotedblleft{}to write, to draft (a text)\textquotedblright{}.

\begin{quote}
Le journaliste a rédigé son article en une heure.
\end{quote}

\begin{quote}
\emph{The journalist wrote his article in an hour.}
\end{quote}

For a text with a purpose: \emph{rédiger une lettre, un rapport}. In daily speech, \emph{écrire}.

\begin{cousinweb}
\emph{rédiger} is a learned borrowing from Latin \emph{redigere} \textquotedblleft{}to bring back, to reduce to order\textquotedblright{} (EtymDB), from \emph{red-} \textquotedblleft{}back\textquotedblright{} and \emph{agere} \textquotedblleft{}to drive\textquotedblright{} (Lewis \& Short).
\end{cousinweb}

\begin{grammarlens}[title={What you've built}]
One more word for telling, explaining and describing at B1.
\end{grammarlens}

\subsection*{Your turn}
\begin{itemize}
\raggedright
  \item \emph{Say it:} \emph{rédiger}
  \item \emph{Say it:} \emph{rédiger}, once more
  \item \emph{Say it:} a sentence of your own with \emph{rédiger}
  \item \emph{Recall:} say the French for the weekly, then the French for the subscriber, then say \emph{rédiger} again
\end{itemize}

\begin{tcolorbox}[breakable,colback=teal!4,colframe=teal!35!black,title={Before you move on}]
What does rédiger mean? (\textquotedblleft{}To write, to draft (a text)\textquotedblright{}.) Say it once more.
\end{tcolorbox}
\section[la rédaction]{la rédaction — the newsroom, the editorial staff; a school essay}

\label{lesson:FR-C279-redaction}



\begin{quote}
Before the new one: say the French for the subscriber, then the French for to write.
\end{quote}

\subsection*{You'll want to know: la rédaction}
\textbf{la rédaction} — \textquotedblleft{}the newsroom, the editorial staff; a school essay\textquotedblright{}.

\begin{quote}
Toute la rédaction a travaillé tard ce soir-là.
\end{quote}

\begin{quote}
\emph{The whole newsroom worked late that evening.}
\end{quote}

At school, \emph{une rédaction} is an essay.

\begin{cousinweb}
\emph{rédaction} is a learned borrowing from Latin \emph{redāctiō} (EtymDB), from \emph{redigere} \textquotedblleft{}to put in order\textquotedblright{}: first the writing of a text, then the people who write it.
\end{cousinweb}

\begin{grammarlens}[title={What you've built}]
One more word for telling, explaining and describing at B1.
\end{grammarlens}

\subsection*{Your turn}
\begin{itemize}
\raggedright
  \item \emph{Say it:} \emph{la rédaction}
  \item \emph{Say it:} \emph{la rédaction}, once more
  \item \emph{Say it:} a sentence of your own with \emph{la rédaction}
  \item \emph{Recall:} say the French for the subscriber, then the French for to write, then say \emph{la rédaction} again
\end{itemize}

\begin{tcolorbox}[breakable,colback=teal!4,colframe=teal!35!black,title={Before you move on}]
What does la rédaction mean? (\textquotedblleft{}The newsroom, the editorial staff; a school essay\textquotedblright{}.) Say it once more.
\end{tcolorbox}
\section[imprimer]{imprimer — to print}

\label{lesson:FR-C279-imprimer}



\begin{quote}
Before the new one: say the French for to write, then the French for the newsroom.
\end{quote}

\subsection*{You'll want to know: imprimer}
\textbf{imprimer} — \textquotedblleft{}to print\textquotedblright{}.

\begin{quote}
Le journal est imprimé pendant la nuit.
\end{quote}

\begin{quote}
\emph{The paper is printed during the night.}
\end{quote}

At home too: \emph{imprimer un billet de train}.

\begin{cousinweb}
\emph{imprimer} is a learned borrowing from Latin \emph{imprimere} \textquotedblleft{}to press into\textquotedblright{} (EtymDB), from \emph{in} and \emph{premere} \textquotedblleft{}to press\textquotedblright{} (Lewis \& Short).
\end{cousinweb}

\begin{grammarlens}[title={What you've built}]
One more word for telling, explaining and describing at B1.
\end{grammarlens}

\subsection*{Your turn}
\begin{itemize}
\raggedright
  \item \emph{Say it:} \emph{imprimer}
  \item \emph{Say it:} \emph{imprimer}, once more
  \item \emph{Say it:} a sentence of your own with \emph{imprimer}
  \item \emph{Recall:} say the French for to write, then the French for the newsroom, then say \emph{imprimer} again
\end{itemize}

\begin{tcolorbox}[breakable,colback=teal!4,colframe=teal!35!black,title={Before you move on}]
What does imprimer mean? (\textquotedblleft{}To print\textquotedblright{}.) Say it once more.
\end{tcolorbox}
\section[paraître]{paraître — to come out, to be published; to seem}

\label{lesson:FR-C279-paraitre}



\begin{quote}
Before the new one: say the French for the newsroom, then the French for to print.
\end{quote}

\subsection*{You'll want to know: paraître}
\textbf{paraître} — \textquotedblleft{}to come out, to be published; to seem\textquotedblright{}.

\begin{quote}
Le magazine paraît tous les jeudis.
\end{quote}

\begin{quote}
\emph{The magazine comes out every Thursday.}
\end{quote}

Two senses: to be published, and \emph{il paraît que} — it seems, people say. The current spelling also allows \emph{paraitre}.

\begin{cousinweb}
\emph{paraître} is inherited from Late Latin \emph{pārēscere}, built on Latin \emph{pārēre} \textquotedblleft{}to appear, to come in sight\textquotedblright{} (Lewis \& Short).
\end{cousinweb}

\begin{grammarlens}[title={What you've built}]
One more word for telling, explaining and describing at B1.
\end{grammarlens}

\subsection*{Your turn}
\begin{itemize}
\raggedright
  \item \emph{Say it:} \emph{paraître}
  \item \emph{Say it:} \emph{paraître}, once more
  \item \emph{Say it:} a sentence of your own with \emph{paraître}
  \item \emph{Recall:} say the French for the newsroom, then the French for to print, then say \emph{paraître} again
\end{itemize}

\begin{tcolorbox}[breakable,colback=teal!4,colframe=teal!35!black,title={Before you move on}]
What does paraître mean? (\textquotedblleft{}To come out, to be published; to seem\textquotedblright{}.) Say it once more.
\end{tcolorbox}
\section[la parution]{la parution — the publication, the release (of a book or paper)}

\label{lesson:FR-C279-parution}



\begin{quote}
Before the new one: say the French for to print, then the French for to come out.
\end{quote}

\subsection*{You'll want to know: la parution}
\textbf{la parution} — \textquotedblleft{}the publication, the release (of a book or paper)\textquotedblright{}.

\begin{quote}
Le livre a eu du succès dès sa parution.
\end{quote}

\begin{quote}
\emph{The book was a success as soon as it came out.}
\end{quote}

\emph{La date de parution} — the publication date.

\begin{cousinweb}
\emph{parution} is built on \emph{paru}, the past participle of \emph{paraître}, which comes from Latin \emph{pārēre} \textquotedblleft{}to appear\textquotedblright{} (Lewis \& Short).
\end{cousinweb}

\begin{grammarlens}[title={What you've built}]
One more word for telling, explaining and describing at B1.
\end{grammarlens}

\subsection*{Your turn}
\begin{itemize}
\raggedright
  \item \emph{Say it:} \emph{la parution}
  \item \emph{Say it:} \emph{la parution}, once more
  \item \emph{Say it:} a sentence of your own with \emph{la parution}
  \item \emph{Recall:} say the French for to print, then the French for to come out, then say \emph{la parution} again
\end{itemize}

\begin{tcolorbox}[breakable,colback=teal!4,colframe=teal!35!black,title={Before you move on}]
What does la parution mean? (\textquotedblleft{}The publication, the release (of a book or paper)\textquotedblright{}.) Say it once more.
\end{tcolorbox}
